\section{رادیکال جیکوبسن}
% AMBIGUITY[L15-Q01]: unresolved-a is the provisional working order. B1-12 match A1-12; B13 ends after the theorem statement on A13. Neither variant is confirmed authoritative.

% Source: sources/lectures/15 Jacobson Radical.pdf (unresolved-a) | PDF page: 1 | item: left and right Jacobson radicals
% Parallel source: sources/lectures/15 Jacobson Radical_241111_015647.pdf (unresolved-b) | PDF page: 1
\begin{frame}[t]{رادیکال‌های چپ و راست جیکوبسن}
\label{lec15:radical-definitions}
\begin{definition}
فرض کنید \(R\) یک حلقهٔ دلخواه باشد. تعریف می‌کنیم:
\[
 J_\ell(R)=\bigcap_{\substack{I\\\text{ایده‌آل چپ ماکسیمال}}} I,
 \qquad
 J_r(R)=\bigcap_{\substack{I\\\text{ایده‌آل راست ماکسیمال}}} I.
\]
این دو را به‌ترتیب \textbf{رادیکال چپ} و \textbf{رادیکال راست جیکوبسن} می‌نامیم.
\end{definition}
ترتیب مطالب از نسخهٔ مفصل‌تر، موقتاً انتخاب شده است: \lr{L15-Q01}.
\end{frame}

% Source: sources/lectures/15 Jacobson Radical.pdf (unresolved-a) | PDF page: 1 | item: integers, division ring, matrix-ring question
% AMBIGUITY[L15-Q03]: both subscripts in the source matrix exercise appear ell; exercise wording retained without silently replacing the second by r.
\begin{frame}[t]{مثال‌های نخستین}
\label{lec15:radical-examples}
\begin{example}
چون \(\mathbb Z\) جابجایی است،
\[
 J_\ell(\mathbb Z)=J_r(\mathbb Z)
 =\bigcap_{\substack{p\\\text{اول}}}(p)=(0).
\]
زیرا اگر هر عدد اولِ \(p\) مقسوم‌علیهِ \(x\) باشد، \(x=0\).
\end{example}
\begin{example}
اگر \(D\) حلقهٔ تقسیم باشد، \((0)\) تنها ایده‌آل چپ و راست ماکسیمال آن است؛ پس
\[
 J_\ell(D)=J_r(D)=0.
\]
\end{example}
\textbf{پرسش جزوه:} رادیکال حلقهٔ \(M_n(D)\) چیست؟
\par\alert{زیرنویس تکراری در صورتِ پرسشِ منبع: \lr{L15-Q03}.}
\end{frame}

% Source: sources/lectures/15 Jacobson Radical.pdf (unresolved-a) | PDF page: 2 | item: IM and its submodule property; red marginal quotient-module caution
\begin{frame}[t]{زیرمدولِ حاصل از عمل یک ایده‌آل}
\label{lec15:ideal-times-module}
\begin{definition}
اگر \(M\) یک \(R\)-مدول چپ و \(I\) یک ایده‌آل چپِ \(R\) باشد، می‌نویسیم
\[
 IM=\left\{\sum_{k=1}^{n}r_km_k\;\middle|\;
 n\in\mathbb N,\ r_k\in I,\ m_k\in M\right\}.
\]
\end{definition}
به‌سادگی دیده می‌شود که \(IM\) زیرمدولی از \(M\) است.
\begin{remark}
اگر \(I\) تنها ایده‌آل چپ باشد، \(R/I\) یک \(R\)-مدول چپ است؛ لازم نیست حلقه باشد.
\end{remark}
\end{frame}

% Source: sources/lectures/15 Jacobson Radical.pdf (unresolved-a) | PDF page: 2 | item: primitive ideals are intersections of maximal left ideals
\begin{frame}[t]{ایده‌آل ابتدایی و ایده‌آل‌های چپ ماکسیمال}
\label{lec15:primitive-intersection}
\begin{lemma}
هر ایده‌آل ابتدایی را می‌توان به صورت اشتراک تعدادی ایده‌آل چپ ماکسیمال نوشت.
همچنین هر ایده‌آل چپ ماکسیمال، شامل یک ایده‌آل ابتدایی است.
\end{lemma}
\textbf{اثبات بخش اول.} اگر \(P\) ابتدایی باشد، برای یک مدول سادهٔ \(M\) داریم
\[
 P=\operatorname{Ann}_R(M)
 =\bigcap_{0\ne x\in M}\operatorname{Ann}_R(x).
\]
برای هر \(0\ne x\in M\)،
\[
 M=Rx\cong R/\operatorname{Ann}_R(x),
\]
پس \(\operatorname{Ann}_R(x)\) ایده‌آل چپ ماکسیمال است. \alert{چرا؟}
\end{frame}

% Source: sources/lectures/15 Jacobson Radical.pdf (unresolved-a) | PDF page: 2 | item: primitive ideal contained in a maximal left ideal
\begin{frame}[t]{ایده‌آل ابتداییِ داخل یک ایده‌آل ماکسیمال}
\label{lec15:maximal-contains-primitive}
\textbf{اثبات بخش دوم.} اگر \(I\) ایده‌آل چپ ماکسیمال باشد، \(M=R/I\) یک \(R\)-مدول ساده است.
بنابراین \(\operatorname{Ann}_R(M)\) ایده‌آلی ابتدایی است.
\medskip

از طرفی، اگر \(a\in\operatorname{Ann}_R(R/I)\)، آنگاه
\[
 a(1+I)=I\quad\Longrightarrow\quad a\in I.
\]
پس
\[
 \operatorname{Ann}_R(R/I)\subseteq I.
\]
در نتیجه \(I\) شامل یک ایده‌آل ابتدایی است. \hfill\(\qed\)
\end{frame}

% Source: sources/lectures/15 Jacobson Radical.pdf (unresolved-a) | PDF page: 3 | item: characterization of left radical
\begin{frame}[t]{توصیف رادیکال با مدول‌های ساده}
\label{lec15:radical-annihilators}
\begin{theorem}
برای هر حلقهٔ \(R\)،
\[
 J_\ell(R)
 =\bigcap_{\substack{P\\\text{ابتدایی}}}P
 =\bigcap_{\substack{M\\\text{سادهٔ چپ}}}\operatorname{Ann}_R(M).
\]
\end{theorem}
\textbf{اثبات.} هر ایده‌آل ابتدایی اشتراک ایده‌آل‌های چپ ماکسیمال است؛ پس
\[
 J_\ell(R)\subseteq\bigcap_{\substack{P\\\text{ابتدایی}}}P.
\]
برای هر ایده‌آل چپ ماکسیمال \(I\)، ایده‌آلی ابتدایی \(P_I\subseteq I\) وجود دارد؛ بنابراین
\[
 \bigcap_{\substack{P\\\text{ابتدایی}}}P\subseteq\bigcap_I P_I
 \subseteq\bigcap_I I=J_\ell(R).
\]
برابری دوم از توصیف ایده‌آل‌های ابتدایی به‌وسیلهٔ نابودسازها به دست می‌آید.
\end{frame}

% Source: sources/lectures/15 Jacobson Radical.pdf (unresolved-a) | PDF page: 4 | item: commutative primitive ideals, two-sidedness, matrix question, annihilation of simple modules
\begin{frame}[t]{چند نتیجه از توصیف رادیکال}
\label{lec15:radical-consequences}
\begin{itemize}
\item اگر \(R\) جابجایی و \(M\) ساده باشد، برای هر \(0\ne x\in M\)،
\[
 \operatorname{Ann}_R(M)=\operatorname{Ann}_R(x).
\]
در نتیجه هر ایده‌آل ابتداییِ یک حلقهٔ جابجایی، ماکسیمال است.
\item چون هر \(\operatorname{Ann}_R(M)\) ایده‌آلی دوطرفه است، \(J_\ell(R)\) نیز دوطرفه است.
\item برای هر مدول چپ سادهٔ \(M\)،
\[
 J_\ell(R)M=0.
\]
\end{itemize}
\textbf{یادآوری پرسش جزوه:} اگر \(D\) حلقهٔ تقسیم باشد، \(J_\ell(M_n(D))\) چیست؟
\end{frame}

% Source: sources/lectures/15 Jacobson Radical.pdf (unresolved-a) | PDF pages: 4-5 | item: finite intersection for a left artinian ring
\begin{frame}[t]{رادیکال در حلقهٔ آرتینی چپ}
\label{lec15:finite-intersection}
\begin{theorem}
اگر \(R\) آرتینی چپ باشد، ایده‌آل‌های چپ ماکسیمال \(I_1,\ldots,I_m\) وجود دارند به طوری که
\[
 J_\ell(R)=\bigcap_{i=1}^{m}I_i.
\]
\end{theorem}
\textbf{اثبات.} هر حلقهٔ یکدار، ایده‌آل چپ ماکسیمال دارد. مجموعهٔ اشتراک‌های متناهیِ این ایده‌آل‌ها را در نظر بگیرید:
\[
 \Sigma=\{I_{i_1}\cap\cdots\cap I_{i_k}\mid k\in\mathbb N\}.
\]
از آرتینی بودن، \(\Sigma\) عضوی مینیمال دارد؛ آن را \(K=\bigcap_{i=1}^m I_i\) بنامید.
برای هر ایده‌آل چپ ماکسیمال \(T\)، \(T\cap K\subseteq K\) و \(T\cap K\in\Sigma\).
پس از مینیمال بودن، \(T\cap K=K\) و \(K\subseteq T\). بنابراین \(K=J_\ell(R)\).
\end{frame}

% Source: sources/lectures/15 Jacobson Radical.pdf (unresolved-a) | PDF page: 5 | item: semisimple iff artinian with zero radical; forward direction
\begin{frame}[t]{نیمه‌سادگی و صفر بودن رادیکال}
\label{lec15:semisimple-criterion}
\begin{theorem}
حلقهٔ \(R\) نیمه‌ساده است اگر و تنها اگر آرتینی چپ باشد و
\[J_\ell(R)=0.\]
\end{theorem}
\textbf{جهت اول.} اگر \(R\) نیمه‌ساده باشد،
\[
 {}_RR=\bigoplus_{i=1}^{n}I_i
\]
که \(I_i\)ها ایده‌آل‌های چپ مینیمال‌اند. هر \(I_i\) ساده و آرتینی است؛ پس \(R\) آرتینی چپ است.
\medskip

همچنین از ساختار ودربرن برای حلقهٔ نیمه‌ساده استفاده می‌کنیم:
\[
 R\cong M_{n_1}(D_1)\times\cdots\times M_{n_k}(D_k),
\]
که \(D_i\)ها حلقهٔ تقسیم‌اند.
\end{frame}

% Source: sources/lectures/15 Jacobson Radical.pdf (unresolved-a) | PDF page: 5 | item: radical of product and matrix rings in forward proof
\begin{frame}[t]{رادیکالِ حاصل‌ضرب حلقه‌ها}
\label{lec15:product-radical}
\textbf{یادآوری جزوه.} ایده‌آل‌های چپ ماکسیمالِ \(R_1\times R_2\) به شکل
\[
 R_1\times\mathfrak m_2\quad\text{یا}\quad\mathfrak m_1\times R_2
\]
هستند که \(\mathfrak m_i\) ایده‌آل چپ ماکسیمالِ \(R_i\) است. پس
\[
 J_\ell(R_1\times R_2)=J_\ell(R_1)\times J_\ell(R_2).
\]
در تجزیهٔ ودربرن، جزوه از \(J_\ell(M_{n_i}(D_i))=0\) استفاده می‌کند؛ در نتیجه
\[
 J_\ell(R)\cong
 J_\ell(M_{n_1}(D_1))\times\cdots\times J_\ell(M_{n_k}(D_k))=0.
\]
\end{frame}

% Source: sources/lectures/15 Jacobson Radical.pdf (unresolved-a) | PDF page: 6 | item: converse via diagonal embedding
\begin{frame}[t]{جهت عکسِ معیار نیمه‌سادگی}
\label{lec15:semisimple-converse}
فرض کنید \(R\) آرتینی چپ و \(J_\ell(R)=0\) باشد. از نتیجهٔ قبل،
\[
 J_\ell(R)=\bigcap_{i=1}^{k}I_i=0
\]
برای ایده‌آل‌های چپ ماکسیمال \(I_i\). همریختی زیر را تعریف می‌کنیم:
\[
 \begin{aligned}
 \varphi:R&\longrightarrow\bigoplus_{i=1}^{k}R/I_i,\\
 r&\longmapsto(r+I_1,\ldots,r+I_k).
 \end{aligned}
\]
چون \(\ker\varphi=\bigcap_{i=1}^kI_i=0\)، قضیهٔ اول همریختی می‌دهد
\[
 {}_RR\cong\operatorname{Im}\varphi.
\]
هر \(R/I_i\) ساده است؛ پس جمع مستقیم بالا نیمه‌ساده و زیرمدول \(\operatorname{Im}\varphi\) نیز نیمه‌ساده است. بنابراین \(R\) نیمه‌ساده است.
\end{frame}

% Source: sources/lectures/15 Jacobson Radical.pdf (unresolved-a) | PDF page: 6 | item: simple left artinian ring is a matrix ring
\begin{frame}[t]{حلقهٔ سادهٔ آرتینی چپ}
\label{lec15:simple-artinian-matrix}
\begin{theorem}
اگر \(R\) یک حلقهٔ ساده و آرتینی چپ باشد، حلقه‌های تقسیم \(D_i\) و اعداد \(n_i\) وجود دارند که
\[
 R\cong M_{n_1}(D_1)\times\cdots\times M_{n_k}(D_k).
\]
\end{theorem}
\textbf{اثبات جزوه.} چون \(R\) ساده و \(J_\ell(R)\) ایده‌آلی دوطرفه و سره است،
\[
 J_\ell(R)=0.
\]
پس با آرتینی بودنِ چپ، \(R\) نیمه‌ساده است و از قضیهٔ ودربرن، تجزیهٔ بالا حاصل می‌شود.
\end{frame}

% Source: sources/lectures/15 Jacobson Radical.pdf (unresolved-a) | PDF page: 7 | item: radical of quotient by the radical
\begin{frame}[t]{خارج‌قسمت بر رادیکال}
\label{lec15:radical-quotient}
\begin{corollary}
برای هر حلقهٔ \(R\)،
\[
 J_\ell\bigl(R/J_\ell(R)\bigr)=0.
\]
اگر \(R\) آرتینی چپ باشد، \(R/J_\ell(R)\) نیمه‌ساده است.
\end{corollary}
\textbf{اثبات.} از تناظر ایده‌آل‌های چپ ماکسیمال،
\[
 J_\ell\bigl(R/J_\ell(R)\bigr)
 =\bigcap_{\substack{\mathfrak m\\\text{چپ ماکسیمال}}}\mathfrak m/J_\ell(R)
 =J_\ell(R)/J_\ell(R)=0.
\]
خارج‌قسمتِ حلقهٔ آرتینی چپ، آرتینی چپ است؛ پس معیار نیمه‌سادگی اعمال می‌شود.
\end{frame}

% Source: sources/lectures/15 Jacobson Radical.pdf (unresolved-a) | PDF page: 7 | item: finiteness of maximal ideals in commutative artinian rings
\begin{frame}[t]{تعداد ایده‌آل‌های ماکسیمال در حالت جابجایی}
\label{lec15:finitely-many-maximals}
\begin{theorem}
حلقهٔ جابجایی و آرتینیِ \(R\)، تنها تعداد متناهی ایده‌آل ماکسیمال دارد.
\end{theorem}
\textbf{اثبات جزوه.} برای ایده‌آل‌های ماکسیمال متمایز \(I_1,I_2,\ldots\)، زنجیر
\[
 I_1\supseteq I_1I_2\supseteq I_1I_2I_3\supseteq\cdots
\]
متوقف می‌شود؛ پس برای یک \(k\)،
\[
 I_1\cdots I_k=I_1\cdots I_k I_{k+1}\subseteq I_{k+1}.
\]
چون \(R\) جابجایی و یکدار است، هر ایده‌آل ماکسیمال اول است؛ بنابراین
\[
 \exists m\leq k:\ I_m\subseteq I_{k+1}.
\]
ماکسیمال بودن، \(I_m=I_{k+1}\) را می‌دهد که با تمایز آن‌ها تناقض دارد.
\end{frame}

% Source: sources/lectures/15 Jacobson Radical.pdf (unresolved-a) | PDF page: 8 | item: first equivalent characterizations of radical
\begin{frame}[t]{توصیف رادیکال با وارون چپ}
\label{lec15:left-inverse-criterion}
\begin{theorem}
اگر \(a\in R\)، موارد زیر معادل‌اند:
\begin{enumerate}
\item \(a\in J_\ell(R)\).
\item برای هر \(r\in R\)، عنصر \(1-ra\) وارون چپ دارد.
\item برای هر \(R\)-مدول چپ سادهٔ \(M\)، \(aM=0\).
\end{enumerate}
\end{theorem}
\textbf{\((1)\Rightarrow(2)\).} اگر \(1-ra\) وارون چپ نداشته باشد، در ایده‌آل چپ ماکسیمالی مانند \(I\) قرار می‌گیرد.
اما \(a\in J_\ell(R)\subseteq I\)؛ پس
\[
 (1-ra)+ra=1\in I,
\]
که تناقض است.
\end{frame}

% Source: sources/lectures/15 Jacobson Radical.pdf (unresolved-a) | PDF page: 8 | item: remaining implications
\begin{frame}[t]{وارون چپ و عمل روی مدول ساده}
\label{lec15:left-inverse-proof}
\textbf{\((2)\Rightarrow(3)\).} اگر \(aM\ne0\)، عنصری \(m\in M\) هست که \(am\ne0\).
از سادگیِ \(M\)، \(M=R(am)\)؛ پس برای یک \(r\in R\)،
\[
 m=ram\quad\Longrightarrow\quad(1-ra)m=0.
\]
چون \(1-ra\) وارون چپ دارد، \(m=0\)، خلاف \(am\ne0\).
\medskip

\textbf{\((3)\Rightarrow(1)\).} برای هر مدول سادهٔ \(M\)،
\[
 aM=0\quad\Longrightarrow\quad a\in\operatorname{Ann}_R(M).
\]
در نتیجه
\[
 a\in\bigcap_{\substack{M\\\text{ساده}}}\operatorname{Ann}_R(M)=J_\ell(R).
\]
\textbf{یادآوری پایین صفحه:} حکم متناظر برای مدول‌های راست و \(J_r(R)\)، با وارون راستِ \(1-ar\)، نیز صحیح است.
\end{frame}

% Source: sources/lectures/15 Jacobson Radical.pdf (unresolved-a) | PDF page: 9 | item: two-sided unit characterization
\begin{frame}[t]{توصیف رادیکال با عناصر وارون‌پذیر}
\label{lec15:unit-criterion}
\begin{theorem}
برای \(a\in R\)،
شرطِ \(a\in J_\ell(R)\) معادل است با این‌که برای هر \(r_1,r_2\in R\)، عنصر
\[1-r_1ar_2\]
وارون‌پذیر باشد.
\end{theorem}
\textbf{جهت بازگشت.} با قرار دادن \(r_2=1\)، معیار قبلی نتیجه می‌دهد.
\medskip

\textbf{جهت رفت.} چون \(J_\ell(R)\) دوطرفه است، \(ar_2\in J_\ell(R)\)؛ پس
\[
 \exists t\in R:\ t(1-r_1ar_2)=1.
\]
بنابراین
\[
 t=1+(tr_1)(ar_2).
\]
\end{frame}

% Source: sources/lectures/15 Jacobson Radical.pdf (unresolved-a) | PDF pages: 9-10 | item: completion of unit proof and equality of radicals
\begin{frame}[t]{برابری رادیکال‌های چپ و راست}
\label{lec15:left-right-equality}
چون \(ar_2\in J_\ell(R)\)، عنصر \(t=1+(tr_1)(ar_2)\) نیز وارون چپ دارد.
با استفاده از \(t(1-r_1ar_2)=1\)، نتیجه می‌شود \(1-r_1ar_2\) وارون‌پذیر است.
\begin{corollary}
\[
 J_\ell(R)=J_r(R).
\]
\end{corollary}
\textbf{اثبات.} معیارِ وارون‌پذیریِ \(1-r_1ar_2\) نسبت به چپ و راست متقارن است؛ پس برای هر دو رادیکال یکسان است.
\medskip

از این پس این ایده‌آل را با \(J(R)\) نشان می‌دهیم و \textbf{رادیکال جیکوبسنِ \(R\)} می‌نامیم؛ نماد \(\operatorname{rad}(R)\) نیز به کار می‌رود.
\end{frame}

% Source: sources/lectures/15 Jacobson Radical.pdf (unresolved-a) | PDF page: 10 | item: nil and nilpotent ideals and ideal powers
\begin{frame}[t]{ایده‌آل پوچ و ایده‌آل پوچ‌توان}
\label{lec15:nil-definitions}
\begin{definition}
ایده‌آل \(I\) از \(R\) را \textbf{پوچ} می‌گوییم هرگاه هر عنصر آن پوچ‌توان باشد:
\[
 \forall a\in I\ \exists n\in\mathbb N:\ a^n=0.
\]
\(I\) را \textbf{پوچ‌توان} می‌گوییم هرگاه برای یک \(n\)، \(I^n=0\).
\end{definition}
\textbf{یادآوری.}
\[
 I^2=\langle ab\mid a,b\in I\rangle
 =\left\{\sum_{\text{متناهی}}a_ib_i\mid a_i,b_i\in I\right\}.
\]
شرط \(I^n=0\) یعنی هر حاصل‌ضربِ \(n\) عنصر از \(I\) صفر باشد.
\begin{remark}
هر ایده‌آل پوچ‌توان، پوچ است؛ عکس آن در حالت کلی صحیح نیست.
\end{remark}
\end{frame}

% Source: sources/lectures/15 Jacobson Radical.pdf (unresolved-a) | PDF pages: 10-11 | item: nil but not nilpotent example
\begin{frame}[t]{ایده‌آلی پوچ که پوچ‌توان نیست}
\label{lec15:nil-not-nilpotent}
\begin{example}
قرار دهید
\[
 A=\mathbb Z[x_1,x_2,x_3,\ldots],\qquad
 I=\langle x_1^2,x_2^3,x_3^4,\ldots\rangle,\qquad R=A/I,
\]
و در \(R\)، ایده‌آل زیر را در نظر بگیرید:
\[
 J=\langle x_1+I,x_2+I,x_3+I,\ldots\rangle.
\]
نمایندهٔ هر عنصرِ \(J\)، چندجمله‌ای با متناهی متغیر است و توانی از این نماینده در \(I\) قرار می‌گیرد؛ پس \(J\) پوچ است.
\medskip

اما \(J\) پوچ‌توان نیست؛ یعنی هیچ \(n\)ای وجود ندارد که \(J^n=0\).
\end{example}
\end{frame}

% Source: sources/lectures/15 Jacobson Radical.pdf (unresolved-a) | PDF page: 11 | item: nil one-sided ideals lie in the radical
\begin{frame}[t]{ایده‌آل‌های پوچ در رادیکال قرار دارند}
\label{lec15:nil-contained-radical}
\begin{theorem}
هر ایده‌آل چپ یا راستِ پوچِ \(R\)، در \(J(R)\) قرار دارد.
\end{theorem}
\textbf{اثبات برای ایده‌آل چپ.} اگر \(a\in I\) و \(r\in R\)، آنگاه \(ra\in I\)؛ پس برای یک \(n\)، \((ra)^n=0\). بنابراین
\[
 1=1-(ra)^n
 =\bigl(1+ra+(ra)^2+\cdots+(ra)^{n-1}\bigr)(1-ra).
\]
پس \(1-ra\) وارون چپ دارد و از معیار قبل، \(a\in J(R)\).
\medskip

حکم ایده‌آل راست به طور متناظر از معیار راست و برابری دو رادیکال به دست می‌آید.
\end{frame}

% Source: sources/lectures/15 Jacobson Radical.pdf (unresolved-a) | PDF pages: 11-12 | item: nilpotence of radical in a left artinian ring, stabilization and Sigma
\begin{frame}[t]{پوچ‌توان بودن رادیکال در حلقهٔ آرتینی چپ}
\label{lec15:artinian-radical-nilpotent}
\begin{theorem}
اگر \(R\) آرتینی چپ باشد، \(J(R)\) پوچ‌توان است.
\end{theorem}
\textbf{اثبات.} به اختصار \(J=J(R)\) بگذارید. از زنجیر نزولی
\[
 J\supseteq J^2\supseteq J^3\supseteq\cdots
\]
و آرتینی بودن، برای یک \(k\)،
\[
 J^k=J^{k+1}=J^{k+2}=\cdots.
\]
ادعا می‌کنیم \(J^k=0\). خلاف آن را فرض کنید و مجموعهٔ زیر را در نظر بگیرید:
\[
 \Sigma=\{I\mid I\text{ ایده‌آل چپِ }R,\ J^kI\ne0\}.
\]
\end{frame}

% Source: sources/lectures/15 Jacobson Radical.pdf (unresolved-a) | PDF page: 12 | item: minimal ideal in the nilpotence proof
\begin{frame}[t]{انتخاب یک ایده‌آل چپ مینیمال}
\label{lec15:radical-minimal-ideal}
چون \(J^kJ=J^{k+1}=J^k\ne0\)، داریم \(J\in\Sigma\)؛ پس \(\Sigma\ne\varnothing\).
از آرتینی بودن، \(\Sigma\) عضوی مینیمال مانند \(I_0\) دارد.
\medskip

چون \(J^kI_0\ne0\)، عنصری \(a_0\in I_0\) وجود دارد که
\[
 J^ka_0\ne0.
\]
از طرفی \(Ja_0\subseteq I_0\) و
\[
 J^k(Ja_0)=J^{k+1}a_0=J^ka_0\ne0.
\]
بنابراین \(Ja_0\in\Sigma\) و از مینیمال بودنِ \(I_0\)،
\[
 Ja_0=I_0.
\]
\end{frame}

% Source: sources/lectures/15 Jacobson Radical.pdf (unresolved-a) | PDF page: 12 | item: contradiction and nil ideals corollary
\begin{frame}[t]{پایان اثبات پوچ‌توان بودن رادیکال}
\label{lec15:radical-nilpotent-conclusion}
چون \(a_0\in I_0=Ja_0\)، عنصری \(t\in J\) وجود دارد که
\[
 ta_0=a_0\quad\Longrightarrow\quad(1-t)a_0=0.
\]
از \(t\in J(R)\)، عنصر \(1-t\) وارون‌پذیر است؛ پس \(a_0=0\)، خلاف \(J^ka_0\ne0\).
بنابراین \(J^k=0\). \hfill\(\qed\)
\begin{corollary}
در هر حلقهٔ آرتینی چپ، هر ایده‌آل چپِ پوچ، پوچ‌توان است.
\end{corollary}
\textbf{دلیل.} اگر \(I\) پوچ باشد، \(I\subseteq J(R)\)؛ چون \(J(R)\) پوچ‌توان است، \(I\) نیز پوچ‌توان است. \alert{چرا؟}
\end{frame}

% Source: sources/lectures/15 Jacobson Radical.pdf (unresolved-a) | PDF page: 13 | item: general Hopkins-Levitzki theorem statement
% Parallel source: sources/lectures/15 Jacobson Radical_241111_015647.pdf (unresolved-b) | PDF page: 13 | item: statement only, no proof
\begin{frame}[t]{قضیهٔ هاپکینز--لویتسکی}
\label{lec15:hopkins-levitzki}
\begin{theorem}
فرض کنید \(J(R)\) پوچ‌توان و \(R/J(R)\) نیمه‌ساده باشد. برای یک \(R\)-مدول چپ \(M\)، موارد زیر معادل‌اند:
\begin{enumerate}
\item \(M\) نوتری است.
\item \(M\) آرتینی است.
\item \(M\) طول متناهی دارد.
\end{enumerate}
\end{theorem}
\textbf{آغاز اثبات.} قبلاً دیده‌ایم که طول متناهی، نوتری و آرتینی بودن را نتیجه می‌دهد؛ پس کافی است هر یک از \((1)\) و \((2)\) را به \((3)\) برسانیم.
\end{frame}

% Source: sources/lectures/15 Jacobson Radical.pdf (unresolved-a) | PDF page: 13 | item: radical filtration and quotient action
\begin{frame}[t]{فیلتراسیونِ مدول با توان‌های رادیکال}
\label{lec15:radical-filtration}
به اختصار \(J=J(R)\) بگذارید و \(n\) را چنان بگیرید که \(J^n=0\). آنگاه
\[
 M\supseteq JM\supseteq J^2M\supseteq\cdots\supseteq J^nM=0.
\]
\textbf{یادآوری جزوه.} اگر \(I\subseteq\operatorname{Ann}_R(N)\)، مدول \(N\) یک \(R/I\)-مدول نیز هست.
\medskip

برای هر \(0\leq k<n\)، قرار دهید
\[
 N_k=J^kM/J^{k+1}M.
\]
چون \(J(J^kM)=J^{k+1}M\)، داریم \(JN_k=0\)؛ پس \(N_k\) یک \(R/J\)-مدول است.
زیرمدول‌های آن بر \(R\) و بر \(R/J\) یکسان‌اند.
\end{frame}

% Source: sources/lectures/15 Jacobson Radical.pdf (unresolved-a) | PDF pages: 13-14 | item: semisimple factors and finite direct sum under the noetherian hypothesis
\begin{frame}[t]{عامل‌های نیمه‌سادهٔ فیلتراسیون}
\label{lec15:filtration-factors}
چون \(R/J\) نیمه‌ساده است، هر \(N_k\) نیز به عنوان \(R/J\)-مدول نیمه‌ساده است؛ بنابراین
\[
 N_k=\bigoplus_{s\in S}N_s,
\]
که \(N_s\)ها مدول‌های سادهٔ ناصفرند.
\medskip

اگر \(M\) نوتری باشد، زیرمدول \(J^kM\) و خارج‌قسمت \(N_k\) نوتری‌اند.
پس در تجزیهٔ بالا، \(S\) متناهی است:
\[
 |S|<\infty.
\]
در نتیجه \(N_k\) دارای سری ترکیبی و طول متناهی است.
\end{frame}

% Source: sources/lectures/15 Jacobson Radical.pdf (unresolved-a) | PDF page: 14 | item: lifting composition series and analogous artinian argument
% AMBIGUITY[L15-Q02]: source writes 0=J^kM/J^{k+1}M and lower quotient M_t/J^kM. Both exponents are provisionally normalized to k+1; see question evidence.
\begin{frame}[t]{از سری عامل‌ها به طول متناهی مدول}
\label{lec15:lift-composition-series}
\alert{خوانش موقتِ توان‌ها در دو فرمولِ جزوه: \lr{L15-Q02}.}
\medskip

سری ترکیبیِ \(N_k\) را به صورت زیر می‌نویسیم:
\[
 0=\frac{J^{k+1}M}{J^{k+1}M}
 \subseteq\frac{M_1}{J^{k+1}M}\subseteq\cdots
 \subseteq\frac{J^kM}{J^{k+1}M}.
\]
از قضیهٔ سوم همریختی، عامل‌های آن به شکل
\[
 \frac{M_{t+1}/J^{k+1}M}{M_t/J^{k+1}M}
 \cong M_{t+1}/M_t
\]
ساده‌اند. با انجام این کار برای هر \(k\)، یک سری ترکیبی برای \(M\) به دست می‌آید؛ پس \(M\) طول متناهی دارد.
\medskip

با استدلال مشابه برای \(M\) آرتینی نیز همین نتیجه به دست می‌آید. \hfill\(\qed\)
\end{frame}

% Source: sources/lectures/15 Jacobson Radical.pdf (unresolved-a) | PDF pages: 14-15 | item: Hopkins theorem for rings and proof
\begin{frame}[t]{قضیهٔ هاپکینز برای حلقه‌ها}
\label{lec15:artinian-implies-noetherian}
\begin{theorem}[هاپکینز، \lr{1939}]
هر حلقهٔ آرتینی چپ، نوتری چپ است.
\end{theorem}
\textbf{اثبات.} اگر \(R\) آرتینی چپ باشد، \(J(R)\) پوچ‌توان است. همچنین
\[
 J\bigl(R/J(R)\bigr)=0
\]
و \(R/J(R)\) آرتینی چپ است؛ بنابراین \(R/J(R)\) نیمه‌ساده است.
\medskip

پس هر دو فرض قضیهٔ قبل برقرارند. آن را برای \(M={}_RR\) به کار می‌بریم؛ چون این مدول آرتینی است، نوتری نیز خواهد بود.
\end{frame}

% Source: sources/lectures/15 Jacobson Radical.pdf (unresolved-a) | PDF page: 15 | item: finitely generated modules over a left artinian ring
\begin{frame}[t]{مدول‌های با تولید متناهی روی حلقهٔ آرتینی}
\label{lec15:fg-modules-artinian}
\begin{corollary}
اگر \(R\) آرتینی چپ باشد، هر \(R\)-مدول چپِ با تولید متناهی، هم نوتری و هم آرتینی است.
\end{corollary}
\textbf{اثبات.} اگر \(M=\sum_{i=1}^{n}Rx_i\)، هر \(Rx_i\) به عنوان \(R\)-مدول، خارج‌قسمتی از \(R\) است:
\[
 Rx_i\cong R/\operatorname{Ann}_R(x_i).
\]
\alert{چرا؟}
پس هر \(Rx_i\) آرتینی و در نتیجه جمع متناهی آن‌ها، یعنی \(M\)، آرتینی است.
از قضیهٔ هاپکینز--لویتسکی، \(M\) نوتری نیز هست.
\begin{corollary}
هر حلقهٔ جابجایی و یکدارِ آرتینی، نوتری است.
\end{corollary}
\end{frame}

% Source: sources/lectures/15 Jacobson Radical.pdf (unresolved-a) | PDF pages: 15-16 | item: nonunital zero-product Prüfer ring counterexample
\begin{frame}[t]{ضرورتِ فرض یکدار بودن}
\label{lec15:nonunital-counterexample}
\begin{remark}
اگر حلقه یکدار نباشد، لزوماً آرتینی بودن، نوتری بودن را نتیجه نمی‌دهد.
\end{remark}
\textbf{مثال جزوه.} گروه \(\mathbb Z_{p^\infty}\) را با جمع گروهی و ضرب صفر در نظر بگیرید:
\[
 xy=0\qquad(x,y\in\mathbb Z_{p^\infty}).
\]
این یک حلقهٔ جابجایی است که آرتینی است، اما نوتری نیست.
\end{frame}

% Source: sources/lectures/15 Jacobson Radical.pdf (unresolved-a) | PDF page: 16 | item: semisimplicity in absence of nonzero nilpotent two-sided ideals
\begin{frame}[t]{نیمه‌سادگی و ایده‌آل‌های دوطرفهٔ پوچ‌توان}
\label{lec15:no-nilpotent-ideals}
\begin{remark}
اگر \(R\) آرتینی چپ باشد و هیچ ایده‌آل دوطرفهٔ ناصفر و پوچ‌توانی نداشته باشد، آنگاه \(R\) نیمه‌ساده است.
\end{remark}
\textbf{دلیل.} چون \(R\) آرتینی چپ است، \(J(R)\) یک ایده‌آل دوطرفهٔ پوچ‌توان است.
از فرض،
\[
 J(R)=0.
\]
در نتیجه از آرتینی بودنِ چپ و صفر بودن رادیکال، \(R\) نیمه‌ساده است.
\end{frame}
